%% secoes/05-conclusao.tex — incluído por main.tex

%%=====================================================================%%
\section{Conclusão}\label{sec:conclusao}
%%=====================================================================%%

Este trabalho aplicou o algoritmo de descoberta de subgrupos SSDP+ a mais de 5 milhões de nascimentos da Coorte de 100 Milhões de Brasileiros vinculada ao SINASC, para quatro desfechos neonatais adversos e oito populações, em 32 análises independentes. 
A análise conjunta mostrou que alguns fatores se repetem entre os desfechos e outros são próprios de cada um. O número reduzido de consultas de pré-natal está associado aos maiores riscos de baixo peso ao nascer, nascimento pré-termo e Apgar baixo, e a primeira gestação acompanha esses três desfechos. A malformação congênita segue padrão distinto, com subgrupos marcados por parto cesáreo, sexo masculino e condições domiciliares presentes, o que sugere influência da detecção e do registro do desfecho.

A análise por segmento mostrou que os grupos vulneráveis não se comportam da mesma forma. As mães menores de 15 anos concentram as maiores prevalências e alguns dos maiores riscos relativos, enquanto na população indígena o nascimento domiciliar se destaca como item de risco para o pré-termo e para o Apgar baixo. Esses padrões não aparecem quando a população é analisada como um todo, o que justifica a análise separada de cada segmento. Os subgrupos encontrados foram disponibilizados em um painel interativo, que permite consultar e comparar os resultados por desfecho e por recorte.
O painel online também incluiu recortes por cidade com o objetivo de contribuir para a tomada de decisão no Brasil de forma mais capilarizada. 

Uma limitação importante deste trabalho é que os dados utilizados vão até 2015, o que pode reduzir a relevância para tomadas de decisões atuais.
Outra limitação é que o estudo não inclui relações de causa e efeito.
Como trabalhos futuros, propõe-se aplicar o método a dados mais recentes e investigar, com métodos adequados de inferência causal, associações relevantes deste estudo, como o papel do pré-natal incompleto e do nascimento domiciliar em populações indígenas.
